\section{Datasets}
\label{sec:datasets}

We use two public corpora with deliberately opposite collection structure. They
are processed independently and never merged; \S\ref{sec:threats} explains why
merging them was itself a source of leakage in an earlier version of this work.

\subsection{5G Traffic}

Seventy-five Wireshark packet captures across fifteen applications in six
service categories, recorded on a commercial 5G network. One capture contains
one application. We read all 330 million packet rows directly from the
distributed archive and extract \emph{conversation segments}: packets between
the capture's client host and one server over one transport, bounded by a
10-second idle timeout and a 10-second active timeout. This yields
164{,}348 flows.

The service label is the dataset's own directory taxonomy
(\texttt{Game\_Streaming}, \texttt{Online\_Game}, \texttt{Live\_Streaming},
\texttt{Stored\_Streaming}, \texttt{Metaverse}, \texttt{Video\_Conferencing}),
so no hand-written mapping intervenes between the collectors' intent and our
labels.

\paragraph{The structural property that matters.} Several applications have only
one or two capture files. Grouping by capture therefore places an entire
application on one side of the split: in our seed-42 session-disjoint partition
the whole of \texttt{roblox} lands in test while training sees only
\texttt{zepeto} for the \texttt{metaverse} class. The 5G session-disjoint
number is consequently a hybrid of session- and application-disjointness. We
report it as such rather than presenting it as a clean session-disjoint result.

\subsection{CESNET-QUIC22}

Twenty-eight daily flow files spanning four consecutive weeks of QUIC traffic
from a national research backbone~\cite{luxemburk2023cesnet}. Each record is
already a bidirectional flow with the first thirty packets as
(inter-arrival, direction, size) triples, plus addresses, ports, destination AS
and the SNI string. We scan all 153 million rows and retain a seeded reservoir
of 400 flows per service category per day, giving 201{,}600 flows across
eighteen categories, perfectly balanced by construction.

\paragraph{Labels come from SNI.} The publishers derive the application and
category fields from the TLS/QUIC server name. Our task is therefore
\emph{predict what SNI would have said, without seeing SNI} --- the right task
for a world with Encrypted Client Hello~\cite{rescorla2024ech}, and one whose
ceiling is set by how well SNI itself categorises traffic. We quantify that
ceiling directly in \S\ref{sec:protocols}.

\subsection{Schema and what the model may see}

Both corpora are reduced to one schema separating \emph{behaviour} from
\emph{provenance}. Behaviour --- per-packet sizes, directions and inter-arrival
times, plus flow-level aggregates --- is what a model may consume. Provenance
--- flow identifier, capture identifier, timestamp, client and server address,
destination port, SNI, AS number, transport protocol --- is retained but
excluded from model inputs, and exists so that the split protocols can group on
it and the probes of \S\ref{sec:protocols} can measure it.

The transport protocol number is provenance rather than input, which is not
obvious and is worth stating. In an earlier merged version of these corpora,
\texttt{protocol} took the value 0 on 11{,}985 rows of which 99.9\% belonged to
one class --- because 0 meant ``the preprocessing script could not parse this
field'', and only one script produced it. The column recorded which pipeline
had run. We now treat it as provenance and probe it like any other identifier.

A test suite asserts mechanically that no provenance column reaches a feature
extractor, and that every split is a partition with no identifier spanning
train and test. It runs in continuous integration and fails the build.
